\begin{abstract}
We characterize the minimax squared-error risk for estimating the observed optimal treatment value from independent and identically distributed observational data with categorical covariates, binary treatments, and binary outcomes. For sample size \(n\), number of covariate cells \(d\ge2\), and public treatment-overlap floor \(0<\epsilon\le1/2\), the risk is of order
\[
\min\left\{1,\frac{d}{n\epsilon\log(ed)}\right\},
\]
with universal comparison constants. The characterization is uniform over unknown cell masses, propensities satisfying the public floor, and conditional outcome means, including treatment-effect ties and null cells. We construct a deterministic observed-data estimator that attains this rate through armwise localization, Jackson polynomial approximation, factorial statistics, and conditional averaging of an auxiliary capped-Poisson construction. The matching lower bound uses weighted approximation and moment-matched probability priors with prescribed common mean, together with a comparison of the full observed likelihood. Consequently, uniform mean-square consistency along growing-cell and shrinking-overlap sequences is possible exactly when \(d/[n\epsilon\log(ed)]\to0\). Under consistency and armwise conditional exchangeability, the same risk characterization applies to the causal oracle value.
\end{abstract}

\section{Introduction}

The population value of individualized treatment choice averages the best conditional outcome attainable in each covariate cell. We establish its minimax squared-error risk in a finite observational experiment with independent and identically distributed observations, categorical covariates, binary treatments, and binary outcomes. With \(n\) the sample size, \(d\ge2\) the number of covariate cells, and \(0<\epsilon\le1/2\) the public treatment-overlap floor, \cref{obj:thm:matched-frontier} gives the rate
\[
\min\left\{1,\frac{d}{n\epsilon\log(ed)}\right\}
\]
with comparison constants universal in all three parameters. The result characterizes the joint cost of an expanding covariate alphabet and weakening overlap for estimation of a scalar optimized population value.

This target follows the optimal-regime and welfare-based treatment-choice formulations of \citet{Murphy2003,Manski2004}. In each cell, it selects the larger of the two conditional outcome means and weights that choice by the population cell probability, as specified in \cref{obj:def:observed-value}. Under consistency, armwise conditional exchangeability, and the public overlap restriction, \cref{obj:prop:identification} identifies this observed functional with the causal oracle value. It also establishes that every admissible observed law has a causal completion. These properties connect the observed-law minimax problem to estimation of the causal oracle value over the corresponding full-data class.

The uniformity domain is central to the result. The class in \cref{obj:def:observed-class} permits arbitrary unknown cell probabilities and conditional outcome means, together with heterogeneous unknown treatment propensities satisfying the public floor in every occupied cell. It includes null cells, boundary outcome means, exact treatment-effect ties, and arbitrarily small positive treatment contrasts. The risk comparison holds for every admissible sample size, alphabet size, and overlap floor. In particular, its constants remain universal as overlap shrinks.

The rate yields a necessary and sufficient condition for uniform mean-square consistency. For sequences \(d_n\), the cell counts, and \(\epsilon_n\), the public overlap floors, satisfying \(d_n\ge2\) and \(0<\epsilon_n\le1/2\), \cref{obj:thm:consistency-threshold} establishes that uniform consistency is possible exactly when
\[
\frac{d_n}{n\epsilon_n\log(e d_n)}\longrightarrow0.
\]
The quantity \(n\epsilon_n\) is the rare-arm sample-size scale allowed by the overlap restriction. The criterion therefore compares that scale with the number of cells divided by its logarithm. For a fixed alphabet, the risk has order \(\min\{1,1/(n\epsilon)\}\), with constants depending on the alphabet; for fixed overlap, it has order \(\min\{1,d/[n\log(ed)]\}\), with constants depending on the overlap floor.

The upper bound is attained by the deterministic armwise factorial estimator in \cref{obj:def:armwise-estimator}. Its active branch uses pilot counts to localize the four observed atom probabilities in each cell, approximates the cell contribution with a tensor Jackson polynomial, and estimates polynomial terms through factorial count statistics. Armwise clipping accounts for the sensitivity of the unknown treatment denominators. Conditional averaging over an auxiliary capped-Poisson count and its allocation marks produces a deterministic function of the observed sample. The parameterized definition also includes empirical estimation for bounded alphabets and a midpoint statistic in the saturated regime. \Cref{obj:thm:observable-upper} guarantees a universal bounded-alphabet cutoff \(D_0\ge2\); the attaining estimator uses the branches specified in \cref{obj:def:armwise-estimator}.

The converse combines weighted polynomial approximation with constrained moment priors. For \(K\), the polynomial degree, and \(M\), the scalar intensity support endpoint, \cref{obj:thm:weighted-separation-and-frontier} establishes that weighted approximation error and target separation between moment-matched probability priors of common mean one have order \(\sqrt M/K\), uniformly for \(K\ge2\), \(2\le M\le K^2\), and \(0\le\epsilon\le1/2\). A localized Jackson packet and common-atom completion construct the required priors. The statistical reduction retains normalization of the sparse alternatives and the information carried by both treatment arms, yielding the lower bound for every measurable estimator in \cref{obj:thm:all-estimator-lower}.

The closest estimand-specific literature develops inference for optimal treatment values under differentiability, convergence, variance, or smoothing conditions \citep[Theorems~1--2]{LuedtkeVanderLaan2016OptimalValue}; \citep[Theorems~2--3]{Shi2020}; \citep[Theorem~3.3 and Corollary~3.4]{WhitehouseChenAusternSyrgkanis2025}. \citet[Theorem~3]{CausalSmithFixedOverlap2026} prove the same-target fixed-overlap rate \(\min\{1,d/[n\log(ed)]\}\) for every fixed \(\epsilon\in(0,1/2)\), with constants allowed to depend on that floor. Our result is uniform for \(0<\epsilon\le1/2\) as the floor shrinks: it identifies the extra \(\epsilon^{-1}\) factor with universal comparison constants. The proof of this extension uses armwise sensitivity control, weighted common-mean priors, and comparison of the full boundary-propensity likelihood. The discrete causal analysis of \citet[Theorems~1--2]{ZengBalakrishnanHanKennedy2024Discrete} provides a nearby comparison for linear arm means and treatment effects. Here the cellwise maximum introduces a nonlinear target. Methodologically, the construction adapts polynomial approximation and unbiased polynomial estimation from discrete-functional estimation \citep{JiaoVenkatHanWeissman2015Functionals,JiaoHanWeissman2018L1} to maxima of unknown observational arm ratios.

\Cref{sec:appendix-lower-bound} develops the statistical lower bound containing \cref{obj:thm:all-estimator-lower}.

The paper proceeds through related work, the observational model and causal interpretation, the minimax rates and consistency criterion, the attaining estimator, and discussion of implications and open questions. The appendices develop identification and the upper bound, weighted approximation and constrained moment priors, the statistical lower bound and verification scope, and proofs of the main results.
